\section{A classical point-packing benchmark}\label{app:point-packing}
The following four values were conjectured by \citet[Section 4]{Anstreicher2009}
and proved by \citet[Propositions 1(i--ii) and 3]{Khajavirad2024}.
We give a variance proof and explicit feasible covariances to show how
an exact relaxation value differs from a spatial certificate bound.
For $n\ge2$ points $(x_i,y_i)\in[0,1]^2$, maximize $\gamma$ subject to
$(x_i-x_j)^2+(y_i-y_j)^2\ge\gamma$ for $i<j$. Replace the two sets of
products by symmetric matrices $X,Y$ and write
\[
 D_{ij}=X_{ii}-2X_{ij}+X_{jj}+Y_{ii}-2Y_{ij}+Y_{jj}\ge\gamma.
\]
RLT imposes all four products of coordinate-bound slacks, including
repeated indices, as in \eqref{eq:box-rlt}. SDP instead imposes
$X\succeq xx^T$, $Y\succeq yy^T$ and the diagonal secants
$X_{ii}\le(\ell_i+u_i)x_i-\ell_i u_i$, and likewise for $Y$.
The symmetry bounds, denoted SYM, restrict $x_i\in[1/2,1]$ for
$i\le n_x=\lceil n/2\rceil$ and $y_i\in[1/2,1]$ for
$i\le k=\lceil n/4\rceil$. Reflection gives the lower-half convention
in Khajavirad's paper. Both the RLT products and the SDP secants are
recomputed on these intervals; adding only first-moment bounds would
define different models.

\begin{proposition}[Four exact values]\label{prop:packing-values}
The optimal values are
\[
\begin{array}{c|c|c}
\text{model}&\text{value}&\text{range}\\\hline
\mathrm{RLT}&2&n\ge2\\
\mathrm{SDP}\text{ (with or without RLT)}&n/(n-1)&n\ge2\\
\mathrm{RLT+SYM}&1/2&n\ge5\\
\mathrm{SDP+SYM}\text{ (with or without RLT)}&k/[4(k-1)]&n\ge5.
\end{array}
\]
Here $k-1=\lfloor(n-1)/4\rfloor$.
\end{proposition}
\begin{proof}
For two coordinates in a common interval $[\ell,\ell+L]$, the RLT
maximum of $X_{ii}-2X_{ij}+X_{jj}$ at fixed means is
\[
 L\min\{x_i+x_j-2\ell,\ 2\ell+2L-x_i-x_j\}\le L^2.
\]
Take both diagonals at their secants and the off-diagonal at the larger
of its two lower McCormick planes. This choice is feasible for all pairs at once: the diagonal secant dominates both
diagonal lower planes, and each pair's lower planes lie below its upper
planes. Thus the bound two holds without symmetry and is attained at
$x=y=\boldsymbol1/2$, with diagonal $1/2$ and off-diagonal zero in each
matrix.
With symmetry, two of the first $k\ge2$ points give the bound $1/2$.
Put the three groups $Q=[k]$, $H=[n_x]\setminus Q$, and
$F=[n]\setminus[n_x]$ at $(3/4,3/4)$, $(3/4,1/2)$, and $(1/2,1/2)$.
Set every diagonal to its interval secant and each off-diagonal to its
lower RLT bound. In group-pair order $QQ,QH,QF,HH,HF,FF$, the resulting
squared distances are
\[
 1/2,\quad7/8,\quad5/4,\quad5/4,\quad13/8,\quad2.
\]
Pairs from groups with fewer than two members impose no extra requirement.
This proves both RLT values.

For SDP, restrict to any $k\ge2$ points sharing $[\ell,\ell+L]$ in one
coordinate, and let $e=\boldsymbol1\in\R^k$. With
$s=\sum_i(x_i-\ell)$, PSD and the secants give
\begin{align*}
 \sum_{i<j}(X_{ii}-2X_{ij}+X_{jj})
 &=k\operatorname{tr}X-e^TXe\\
 &\le k[(2\ell+L)\sum_i x_i-k\ell(\ell+L)]-(\sum_i x_i)^2\\
 &=kLs-s^2\le k^2L^2/4.
\end{align*}
The pair average per coordinate is at most $kL^2/[2(k-1)]$.
Adding the two coordinates, first with $k=n,L=1$ and then for the
symmetry subset with $L=1/2$, gives the two asserted SDP upper bounds.

Without symmetry, set $x=y=e/2$ and
\begin{equation}\label{eq:packing-covariance}
 X=Y=ee^T/4+Z,\qquad
 Z=\frac{n}{4(n-1)}(I-ee^T/n)\succeq0.
\end{equation}
The diagonals equal $1/2$ and the off-diagonals equal
$1/4-1/[4(n-1)]$, so every pair has total distance $n/(n-1)$.
These entries also satisfy all RLT bounds.

With symmetry, use the same three group means as above. In each coordinate
set the moment matrix to mean outer product plus a covariance that is
block diagonal in $Q,H,F$. On $Q$ use in both coordinates
\[
 Z_Q=\frac{k}{16(k-1)}(I-ee^T/k).
\]
On $H$ use $I/16$ in $x$ and $I/4$ in $y$; on $F$ use $I/4$ in both.
These matrices are PSD because $I-ee^T/k$ is an orthogonal projection.
Every diagonal equals its tightened secant. The group-pair distances are
\[
 \frac{k}{4(k-1)},\quad\frac12,\quad\frac34,\quad
 \frac58,\quad\frac78,\quad1.
\]
They are at least the first value for every $k\ge2$.
For RLT feasibility, cross-group and distinct $H,F$ entries are products
of their means and hence obey the four planes. On $Q$ an off-diagonal is
$9/16-1/[16(k-1)]\in[1/2,5/8]$, exactly within the tightened RLT range
at mean $3/4$. Diagonal secants satisfy their diagonal lower planes as
above. Thus adding RLT cannot change either SDP optimum.
\end{proof}

In dimension $d$, the same one-coordinate average and the same covariance
in each coordinate give RLT value $d$ and SDP value
$dn/[2(n-1)]$ on $[0,1]^d$, for $n\ge2$. A subset of $k\ge2$ points
in a common side-$L$ cube, including $L=0$, forces an SDP bound at most
$dL^2k/[2(k-1)]$. These are dimension-scaled averaging statements; they
do not cover other symmetry or ordering models, or symmetry values for
$n<5$. The benchmark explains the weakness of these particular root
relaxations; it gives no lower bound for every possible strengthening
or spatial algorithm.
